% fixture: amsmath/multline-label-tag
% tier: T3
% engine: document
% covers: \label and \tag in multline and multline* (read in \mmeasure@: \label@in@display, \tag@in@align; \make@display@tag in the last row puts the tag and then the \label), a \tag* that fits beside the last row and a \tag that does not (\shifttag@: below it), with \eqref and \ref to them
\documentclass{article}
\usepackage{amsmath}
\begin{document}
A long formula is split with multline; see \eqref{m:one}, \eqref{m:two} and \ref{m:three}.
\begin{multline} \label{m:one}
a + b + c + d + e + f + g + h + i + j + k + l + m + n + o + p \\
= q + r + s + t + u + v + w + x + y + z
\end{multline}
With a tag in place of the number,
\begin{multline}
a + b + c + d + e + f + g + h + i + j + k + l + m + n + o + p \tag{M} \label{m:two} \\
= q + r + s + t + u + v + w + x + y + z
\end{multline}
and in the starred form, with a tag too wide to stand beside the last row:
\begin{multline*}
a + b + c + d + e + f + g + h + i + j + k + l + m + n + o + p \\
= q + r + s + t + u + v + w + x + y + z + a + b + c + d + e + f + g + h + i + j + k + l + m + n + o + p \tag*{[the long one]} \label{m:three}
\end{multline*}
The numbering goes on:
\begin{multline}
a + b + c + d + e + f + g + h + i + j + k + l + m + n + o + p \\
= q + r + s + t + u + v + w + x + y + z
\end{multline}
which ends the example.
\end{document}
